\section{Future Plan}
\label{sec:future}

The reproduction motivates a focused hypothesis: \textbf{jointly evolving comparator guidance and a bounded selection-and-batch policy can improve solution quality per unit of total search time.} Low confidence left all ten candidates tied in three batches, while generating a proposal took longer on average than executing a selected program. The initial study will therefore target the usefulness of comparisons and the cost of obtaining alternatives. The underlying LLMs will remain frozen.

\paragraph{Establishing a common evaluation signal.}
The observed 0.00971 private-error selection gap and inconsistent program-defined validation motivate a shared evaluation layer. All experimental arms will use the same held-out split of public training data and a harness-computed metric matching the benchmark formula, target columns, and scored positions. Candidate programs will receive only the training partition with labels and the held-out features; held-out labels will remain in an evaluator-only location outside the execution sandbox. Programs will supply predictions with checked IDs and shapes, and the evaluator will compute the score. This common utility will replace parsed self-reported metrics in Eqs.~\ref{eq:objective}--\ref{eq:target} for the planned experiments. The unordered-cache bug will also be fixed by storing winner identities and mapping them to the current input order. These shared changes will be documented separately from adaptation and applied to every comparison arm; the original reproduction will remain unchanged.

\paragraph{Evolving comparison and resource allocation.}
The comparator receives neither explicit parent outcomes nor the journal; combined with the low-confidence rate, this motivates a persistent skill summarizing executed revisions, standardized outcomes, and prediction errors. Evolution will occur \emph{within each run}: after a completed improvement batch, one update call may revise the skill and choose $m\in\{2,5,10\}$ and $\tau\in\{0.5,0.7,0.9\}$. Smaller batches can reduce proposal cost when ranking fails to discriminate. The static harness will be revised once from the reproduction's evidence before the pilot and then frozen; the evolving arm will start from this identical version.

A revision will be tried on the next batch without restarting search. Before execution, the revised comparator will also predict whether each selected candidate improves on its valid parent; standardized execution outcomes will then check those predictions. The revision will be retained only if at least one such prediction is testable, all testable preferences are correct, and the increase in running-best standardized validation utility per elapsed batch minute is at least that of the preceding batch; otherwise the prior harness will be restored. Trial time includes update calls, generation, comparisons, execution, and parsing. This is a prespecified local acceptance heuristic; the controlled comparisons below test whether it improves overall search.

The pilot will allow at most two revision attempts, each with a five-minute update-call limit, and initiate none after 2.5 hours. Given the observed proposal cost, the planning estimate is zero to two completed revisions in four hours, depending on when improvement batches become available. An unfinished or untestable trial will leave the prior version active; every attempt, including a rejected one, consumes the same four-hour allowance. The number of completed revisions will be reported as a feasibility outcome.

Rehearse provides the closest reference for retrieving modifications and outcomes, while Evo-Harness motivates maintaining and revising reusable guidance \citep{ji2026rehearse,wei2026evoharness}. AIDE$^2$ supplies a precedent for evolving search policies \citep{srikanth2026aide2}. The proposed emphasis is their interaction with ForeAgent's confidence-gated predictive filter under measured generation and execution costs. More confident judgments alone will not count as improvement: diagnostics will assess agreement with standardized execution outcomes, confident errors, failed executions per batch, and time spent generating candidates. Outcomes for skipped programs will remain unknown unless acquired through an explicitly budgeted verification sample.

\paragraph{Experimental scale and success criteria.}
The initial comparison will use Stanford COVID Vaccine as a development task, three paired local random seeds, and four-hour search windows on the same models, hardware, and API configuration. The four-hour window covers the reproduced validation best at 3.46 hours and lies within the period of comparable search breadth before 4.35 hours (Section~\ref{sec:reproduction}); it targets early search and leaves later behavior for a longer follow-up. Pairing fixes the data split and local random streams for controller choices, rematches, and Bernoulli gating; API-generated LLM outputs remain stochastic, and divergent trajectories can consume those streams differently. Three arms will compare the shared-fix ForeAgent baseline, a static revised harness, and an evolving harness. This specifies 36 hours of primary search windows, including online adaptation and next-batch trials. Initial harness construction and the independent target-setting pilot will be itemized separately and included in aggregate costs. Evolving versus shared-fix ForeAgent directly addresses the research question; evolving versus static guidance isolates the value of continued revision. Comparator-only and batch-policy-only ablations will follow if the main contrast warrants further study.

Primary outcomes will be standardized validation quality at a fixed total budget and the probability of reaching a prespecified target within that budget. For each task, the target will be fixed from an independent baseline pilot's two-hour validation best before comparative runs begin; pilot costs will be reported. Runs that miss the target will be treated as censored at the budget limit. This three-pair study is a pilot. A prespecified positive signal requires lower capped hitting time $\min(T_\pi(\delta),B)$ in all three pairs against the shared-fix baseline, with no worse standardized endpoint error in any pair. Equal capped times count as ties, not improvements. All paired time and error differences and target-hit counts will be reported; this rule identifies a consistent direction for follow-up rather than statistical significance. The evolving-versus-static contrast will use the same descriptive rule. Stanford's early Beat Ratio ceiling makes raw validation error and time-to-target more informative than leaderboard-threshold attainment. Raw private MCRMSE still distinguishes final solutions.

A second-stage transfer evaluation will use one additional single-GPU MLE-bench task excluded from harness development and prior skill construction, selected before comparative testing for tractable training and non-saturated score trajectories. It will compare the static and evolving arms with three paired seeds and the same four-hour limit. Development costs will count toward the aggregate transfer budget, alongside GPU usage and API expenditure. End-to-end efficiency claims will additionally compare arms at matched aggregate budgets, charging harness construction and revision evaluation to the respective arms and allocating the same total allowance to the static arms. Private labels will remain unavailable to harness revision and will be used only for final reporting.

\paragraph{Extensions guided by the main result.}
A matched-input comparison audit can first check whether the low-confidence pattern persists when candidate inputs are held fixed across service configurations; the current archive cannot isolate a provider effect. Subsequent work may extend the persistent skill to generation, testing whether it reduces failed executions per batch without increasing generation time, or evolve parent selection and debugging allocation. These extensions preserve the broader end-to-end search-efficiency question while keeping the first experiment focused on the bottlenecks measured in this reproduction.
